% 中文阅读版的共享排版设置；使用 XeLaTeX 编译。
\usepackage{fontspec}
\usepackage{xeCJK}
\setmainfont{NimbusRoman-Regular.otf}[Path=/usr/share/fonts/opentype/urw-base35/,BoldFont=NimbusRoman-Bold.otf,ItalicFont=NimbusRoman-Italic.otf,BoldItalicFont=NimbusRoman-BoldItalic.otf]
\setsansfont{NimbusSans-Regular.otf}[Path=/usr/share/fonts/opentype/urw-base35/,BoldFont=NimbusSans-Bold.otf,ItalicFont=NimbusSans-Italic.otf,BoldItalicFont=NimbusSans-BoldItalic.otf]
\setmonofont{DejaVu Sans Mono}[Scale=MatchLowercase]
\setCJKmainfont{Noto Serif CJK SC}[BoldFont={Noto Serif CJK SC Bold},ItalicFont={Noto Serif CJK SC}]
\setCJKsansfont{Noto Sans CJK SC}[BoldFont={Noto Sans CJK SC Bold}]
\setCJKmonofont{Noto Sans Mono CJK SC}[BoldFont={Noto Sans Mono CJK SC Bold},Scale=0.9908]
\renewcommand{\abstractname}{摘要}
\renewcommand{\figurename}{图}
\renewcommand{\tablename}{表}
\renewcommand{\appendixname}{附录}
\renewcommand{\refname}{参考文献}
\providecommand{\proofname}{证明}
\renewcommand{\proofname}{证明}
% ACL 样式中的匿名作者栏和参考文献标题采用固定英文文本，在此本地化。
\usepackage{etoolbox}
\makeatletter
\patchcmd{\outauthor}{Anonymous ACL submission}{ACL 匿名投稿}{}{}
\patchcmd{\thebibliography}{References\@mkboth{References}{References}}{参考文献\@mkboth{参考文献}{参考文献}}{}{}
\makeatother
\emergencystretch=1em
